\section*{Sets and Indices}

\[
T=\{1,\dots,8\} \quad \text{weeks}
\]

\[
P=\{I,II\} \quad \text{product families}
\]

Training can start only in weeks
\[
T^{\text{start}}=\{1,\dots,7\},
\]
because the code defines training-start variables only for \(t\le T-1\).

\section*{Parameters}

All parameter values are read from \texttt{general\_parameters.csv} and \texttt{table\_1.csv}.

\[
S_0 \quad \text{(code: \texttt{S0}) = initial skilled worker count}
\]

\[
r_I,\; r_{II} \quad \text{(code: \texttt{rate\_I}, \texttt{rate\_II}) = production rates in kg/hour}
\]

\[
H^N \quad \text{(code: \texttt{normal\_hours}) = regular weekly hours per worker}
\]

\[
H^{OT} \quad \text{(code: \texttt{overtime\_hours}) = overtime-week hours per overtime skilled worker}
\]

\[
c^{tr} \quad \text{(code: \texttt{train\_cap}) = number of trainees produced by one trainer per 2-week class}
\]

\[
L \quad \text{(code: \texttt{train\_period}) = training duration in weeks} = 2
\]

\[
w^S \quad \text{(code: \texttt{wage\_skilled}) = weekly wage of a skilled worker on regular duty or as trainer}
\]

\[
w^{OT} \quad \text{(code: \texttt{wage\_overtime}) = weekly wage of a skilled worker assigned to overtime}
\]

\[
w^{TD} \quad \text{(code: \texttt{wage\_trainee\_during}) = weekly wage of a trainee during training}
\]

\[
w^{TA} \quad \text{(code: \texttt{wage\_trainee\_after}) = weekly wage of a trained worker after graduation}
\]

\[
\gamma_I,\; \gamma_{II} \quad \text{(code: \texttt{comp\_I}, \texttt{comp\_II}) = weekly compensation cost per kg of backlog}
\]

\[
G \quad \text{(code: \texttt{new\_worker\_goal}) = required number of newly trained workers by the end of week 8}
\]

\[
d_{p,t} \quad \text{weekly demand for product } p \text{ in week } t
\]

The implementation also uses a big-\(M\) constant
\[
M = S_0 H^{OT} + S_0 c^{tr} H^N + 10000
\]
(code: \texttt{BIG\_M}).

For compactness, define the code-induced cumulative expressions
\[
B_t^{tr} = \sum_{\substack{s\in T^{\text{start}}:\\ s\le t \le s+L-1}} n_s
\]
as the number of trainers busy in week \(t\), and
\[
G_t^{av} = \sum_{\substack{s\in T^{\text{start}}:\\ s+L \le t}} c^{tr} n_s
\]
as the number of graduates available for production in week \(t\).

Also define the number of trainees currently in training in week \(t\):
\[
Q_t^{tr} = \sum_{\substack{s\in T^{\text{start}}:\\ s\le t \le s+L-1}} c^{tr} n_s.
\]

\section*{Decision Variables}

\[
n_t \in \mathbb{Z}_+ \quad \text{for } t\in T^{\text{start}}
\]
(code: \texttt{n\_train[t]}): number of skilled workers who start a training class in week \(t\).

\[
o_t \in \mathbb{Z}_+ \quad \text{for } t\in T
\]
(code: \texttt{n\_overtime[t]}): number of skilled workers assigned to overtime in week \(t\).

For each \(t\in T\), nonnegative hour-allocation variables:
\[
h^S_{I,t},\, h^S_{II,t} \ge 0
\]
(code: \texttt{h\_I\_s[t]}, \texttt{h\_II\_s[t]}): regular skilled-worker production hours.

\[
h^{OT}_{I,t},\, h^{OT}_{II,t} \ge 0
\]
(code: \texttt{h\_I\_ot[t]}, \texttt{h\_II\_ot[t]}): overtime skilled-worker production hours.

\[
h^G_{I,t},\, h^G_{II,t} \ge 0
\]
(code: \texttt{h\_I\_g[t]}, \texttt{h\_II\_g[t]}): production hours supplied by graduated workers.

\[
b_{I,t},\, b_{II,t} \ge 0 \quad \text{for } t\in T
\]
(code: \texttt{backlog\_I[t]}, \texttt{backlog\_II[t]}): end-of-week backlog for each product.

\[
y_t \in \{0,1\} \quad \text{for } t\in T
\]
(code: \texttt{produce\_I\_flag[t]}): weekly family-selection flag; \(y_t=1\) allows only Food I production, \(y_t=0\) allows only Food II production.

\section*{Objective}

Minimize total 8-week cost:
\[
\min \sum_{t\in T}
\Big[
(w^S)(S_0-B_t^{tr}-o_t)
+ (w^{OT})o_t
+ (w^S)B_t^{tr}
+ (w^{TD})Q_t^{tr}
+ (w^{TA})G_t^{av}
+ \gamma_I b_{I,t}
+ \gamma_{II} b_{II,t}
\Big].
\]

This exactly matches the implemented cost accounting, which separately charges:
regular skilled-worker wages for available non-overtime skilled workers,
overtime wages for overtime workers,
trainer wages for busy trainers,
trainee wages during training,
graduate wages after training,
and weekly compensation on outstanding backlog.

\section*{Constraints}

For all \(t\in T\), overtime cannot exceed available skilled workers:
\[
o_t \le S_0 - B_t^{tr}.
\]

For all \(t\in T\), regular skilled-worker production hours are limited by available non-overtime skilled workers:
\[
h^S_{I,t} + h^S_{II,t}
\le (S_0 - B_t^{tr} - o_t)H^N.
\]

For all \(t\in T\), overtime production hours are limited by the number of overtime workers:
\[
h^{OT}_{I,t} + h^{OT}_{II,t}
\le o_t H^{OT}.
\]

For all \(t\in T\), graduate production hours are limited by available graduates:
\[
h^G_{I,t} + h^G_{II,t}
\le G_t^{av} H^N.
\]

Mutual exclusion of product families in each week, for all \(t\in T\):
\[
h^S_{I,t} + h^{OT}_{I,t} + h^G_{I,t} \le M y_t,
\]
\[
h^S_{II,t} + h^{OT}_{II,t} + h^G_{II,t} \le M(1-y_t).
\]

Thus, in each week, all productive hours from every workforce source are assigned either to Food I only, or to Food II only, or to neither product if all such hours are zero.

For all \(t\in T\), backlog propagation for Food I:
\[
b_{I,t} \ge b_{I,t-1} + d_{I,t} - r_I\left(h^S_{I,t}+h^{OT}_{I,t}+h^G_{I,t}\right),
\]
with \(b_{I,0}=0\).

For all \(t\in T\), backlog propagation for Food II:
\[
b_{II,t} \ge b_{II,t-1} + d_{II,t} - r_{II}\left(h^S_{II,t}+h^{OT}_{II,t}+h^G_{II,t}\right),
\]
with \(b_{II,0}=0\).

For all \(t\in T\), trainer load cannot exceed the skilled workforce:
\[
S_0 - B_t^{tr} \ge 0.
\]

Training-goal requirement by the end of week 8:
\[
\sum_{\substack{s\in T^{\text{start}}:\\ s+L-1\le 8}} c^{tr} n_s \ge G.
\]

Since \(L=2\) and \(T^{\text{start}}=\{1,\dots,7\}\), this sums all training starts \(s=1,\dots,7\), including week 7 starts that finish by the end of week 8 but do not contribute productive hours within the horizon.

\section*{Notes on Implementation}

The source code solves this formulation as a mixed-integer linear program using PuLP-style modeling and \texttt{GUROBI\_CMD}.

The model does \emph{not} include explicit shipment variables. Instead, backlog variables satisfy lower-bound recursions
\[
b_{p,t} \ge b_{p,t-1}+d_{p,t}-\text{production}_{p,t},
\]
and because backlog carries positive cost in the objective, optimal solutions drive these inequalities as tight as beneficial.

The weekly wage terms follow the code literally:
\[
(w^S)(S_0-B_t^{tr}-o_t) + (w^{OT})o_t + (w^S)B_t^{tr}
= (w^S)(S_0-o_t) + (w^{OT})o_t,
\]
so trainers are paid as skilled workers, and workers assigned to overtime receive the overtime weekly wage instead of the regular weekly wage.

Graduates become available only from week \(s+L\) onward after a training start in week \(s\), exactly as implemented in \texttt{graduates\_available(t)}. Trainees in classes started in week 7 therefore count toward the end-of-horizon training goal but do not supply production hours in week 8.

The mutual exclusion revision is implemented with a single binary variable \(y_t\) and big-\(M\) constraints that force all production hours in week \(t\) to belong to only one food family.
